\chapter{Proposta}
\label{cap:proposta}

% -----------------------------------------------------------------
% Objetivo do capítulo:
% Apresentar a proposta da dissertação: primeiro o problema de
% pesquisa formalizado e os objetivos, depois a arquitetura
% concreta da ferramenta. É aqui que entram diagramas de
% arquitetura, pseudocódigo dos indicadores, exemplos de saída,
% decisões de engenharia e ajustes feitos durante o
% desenvolvimento.
%
% Referência interna:
%   docs/pipeline/README.md
%   docs/pipeline/resumo_algoritmos.md
%   docs/pipeline/resumo_metodo_estatistico_global.txt
% -----------------------------------------------------------------


\section{Problema de pesquisa}
\label{sec:problema_pesquisa}

% TODO: enunciar o problema de pesquisa de forma autocontida.
% Cobertura prevista:
% - Mecanismos sobrepostos (materia-prima, equipamento, combinacao) que produzem fenotipos semelhantes na fratura, referenciando a Secao~\ref{sec:problema_rompimentos}.
% - Ausencia de rotulo de causa nos dados operacionais, referenciando a Secao~\ref{sec:dados_disponiveis}.
% - Enunciar a pergunta central:
%   "Como combinar sinais operacionais de maquina e evidencias estatisticas historicas de rompimentos para inferir a provavel origem de eventos de rompimento em uma fabrica de fios?"


\section{Objetivos}
\label{sec:objetivos}


\subsection{Objetivo geral}
\label{subsec:objetivo_geral}

% TODO (versao planejada): Desenvolver uma ferramenta hibrida de apoio a decisao que combine sinais operacionais de maquina e evidencias estatisticas historicas de rompimentos para inferir a provavel origem de eventos de rompimento em uma fabrica de fios esmaltados de cobre e aluminio.


\subsection{Objetivos específicos}
\label{subsec:objetivos_especificos}

% TODO (versao planejada com 5 itens):
% 1. Conduzir revisao sistematica da literatura sobre deteccao de anomalias e analise estatistica aplicadas a atribuicao de causa em processos industriais com ausencia de rotulos.
% 2. Caracterizar e preparar as fontes de dados operacionais e historicos disponiveis para analise.
% 3. Desenvolver um eixo de deteccao de anomalias operacionais a partir de series temporais de variaveis sensoriais.
% 4. Desenvolver um eixo de analise estatistica historica de rompimentos segmentado por lote, material, fornecedor e maquina.
% 5. Integrar os dois eixos em uma ferramenta de inferencia e avaliar o resultado a partir de casos historicos da fabrica.


\section{Visão geral da arquitetura}
\label{sec:arquitetura}

% Conteúdo a desenvolver:
%   - Nome interno da ferramenta (WireSense) e organização em
%     dois pilares (Pilar MP e Pilar EQ).
%   - Diagrama da arquitetura (referenciar imagem em
%     docs/diagramas/).
%   - Papel de cada pilar e como se comunicam.


\section{Pipeline de dados}
\label{sec:pipeline}

% Conteúdo a desenvolver:
%   - Sequência de etapas, das tabelas brutas ao diagnóstico.
%   - Mapa entre scripts implementados e etapas do pipeline.
%   - Pontos de integração com as bases existentes.


\section{Implementação do módulo de anomalia (Pilar EQ)}
\label{sec:impl_eq}

% Conteúdo a desenvolver:
%   - Conjunto de 14 sensores e justificativa técnica da
%     escolha.
%   - Treinamento e ajuste do Isolation Forest.
%   - Critérios de SEM_MODELO e fallback ao Pilar MP.
%   - Exemplos de saída do módulo.


\section{Implementação do módulo estatístico (Pilar MP)}
\label{sec:impl_mp}

% Conteúdo a desenvolver:
%   - Indicadores calculados (taxas, exposição, risco relativo,
%     cruzamento entre máquinas).
%   - Decisão metodológica registrada em 12/05: o Pilar MP
%     conta apenas rompimentos, com bolha removida do
%     numerador.
%   - Regra MODERADA: o Pilar MP é acionável, o ground truth
%     do relatório de devoluções é reservado para validação
%     externa.
%   - Exemplos de saída do módulo.


\section{Regras de integração e diagnóstico final}
\label{sec:regras_integracao}

% Conteúdo a desenvolver:
%   - Lógica de combinação entre Pilar EQ e Pilar MP.
%   - Categorias de origem provável.
%   - Forma de apresentação do diagnóstico (texto, indicadores,
%     evidência de apoio).


\section{Decisões de engenharia e ajustes iterativos}
\label{sec:decisoes_engenharia}

% Conteúdo a desenvolver:
%   - Principais decisões tomadas durante o desenvolvimento.
%   - Compromissos entre cobertura, precisão e
%     interpretabilidade.
%   - Pontos de extensibilidade da ferramenta.
